% Part: second-order-logic
% Chapter: syntax-and-semantics
% Section: terms-formulas

\documentclass[../../../include/open-logic-section]{subfiles}

\begin{document}

\olfileid{sol}{syn}{frm}

\olsection{ٹرم تے \printtoken{P}{formula}}

پہلے درجے دی منطق وانگ، دوجے درجے دی منطق دیاں لکھتاں
بنیادی علامتاں دی اک فہرست توں بنائیاں جاندیاں نیں۔ ایس وچ
\emph{!!{variable}s}، \emph{!!{constant}s}،
\emph{!!{predicate}s} تے کدے کدے
\emph{!!{function}s} شامل نیں۔ انہاں نوں منطقی رابطیاں،
مسوّراں، تے قوساں تے کومیاں ورگیاں لکھت دیاں نشانیاں
نال ملا کے \emph{ٹرم} تے \emph{!!{formula}s}
بنائے جاندے نیں۔ فرق ایہہ اے کہ دوجے درجے دی منطق وچ
شے دے متغیراں توں علاوہ تعلقاں تے فنکشوناں دے
متغیر وی ہوندے نیں، تے انہاں اُتے مسوّریت دی
اجازت وی ہوندی اے۔ سو دوجے درجے دی منطق دیاں منطقی
علامتاں پہلے درجے دی منطق والیاں نیں، نال ایہہ وی:

\begin{enumerate}
\item ہر رُکن-گنتی~$n$ لئی دوجے درجے دے تعلقی !!{variable}s دا
  اک !!{denumerable}s سیٹ:
  $\Obj V_0^n$, $\Obj V_1^n$, $\Obj V_2^n$, \dots
\item دوجے درجے دے فنکشنی !!{variable}s دا اک
  !!{denumerable}s سیٹ:
  $\Obj u_0^n$, $\Obj u_1^n$, $\Obj u_2^n$, \dots
\end{enumerate}

جیویں اسیں $x$, $y$, $z$ نوں پہلے درجے دے
متغیراں $\Obj v_i$ لئی ماورا متغیراں دے طور تے ورتدے آں،
اُویں $X$, $Y$, $Z$ وغیرہ نوں $\Obj V_i^n$ لئی
تے $u$, $v$ وغیرہ نوں~$\Obj u_i^n$ لئی ماورا متغیراں
دے طور تے ورتاں گے۔

\begin{explain}
دوجے درجے دی زبان دیاں غیر منطقی علامتاں اوہو
طریقے نال متعین کیتیاں جاندیاں نیں جیویں پہلے درجے
دی زبان دیاں: اوہدے !!{constant}s، !!{function}s
تے !!{predicate}s دی فہرست دے کے۔

پہلے درجے دی منطق وچ !!{identity}~$\eq$ نوں عموماً
شامل کیتا جاندا اے۔ پہلے درجے دی منطق وچ
زبان~$\Lang{L}$ دیاں غیر منطقی علامتاں کسے وی
دلچسپ گل نوں بیان کرن لئی بہت ضروری نیں۔ بے شک
کجھ !!{sentence}s غیر منطقی علامتاں توں بناں وی
ہوندے نیں، پر صرف~$\eq$ نال کوئی دلچسپ گل کہنا
اوکھا اے۔ دوجے درجے دی منطق وچ، تعلقی تے فنکشنی متغیر بے حد نیں۔ ایس لئی اسیں
وکھریاں غیر منطقی علامتاں توں بناں وی اوہ کجھ آکھ سکدے آں
جیہڑا پہلے درجے دی کسے زبان وچ آکھ سکدے آں۔
\end{explain}

\begin{defn}[دوجے درجے دے ٹرم]
زبان~$\Lang L$ دے \emph{دوجے درجے دے ٹرماں} دا سیٹ،
$\TrmSOL[L]$، \olref[fol][syn][frm]{defn:terms}
دی تعریف وچ ایہہ شق شامل کر کے متعین ہوندا اے:
\begin{enumerate}
\item جے $u$ اک $n$-رکنی فنکشنی متغیر اے تے $t_1$, \dots, $t_n$
  ٹرم نیں، تاں $\Atom{u}{t_1, \ldots, t_n}$ اک ٹرم اے۔
\end{enumerate}
\end{defn}

\begin{explain}
سو دوجے درجے دا ٹرم پہلے درجے دے ٹرم ورگا اے۔
فرق ایہ اے کہ جتھے پہلے درجے دے ٹرم
وچ !!a{function}~$\Obj{f^n_i}$ آؤندا اے، اُتھے
دوجے درجے دے ٹرم وچ اُس دی تھاں فنکشنی
متغیر~$\Obj{u^n_i}$ آ سکدا اے۔
\end{explain}

\begin{defn}[دوجے درجے دے \usetoken{s}{formula}]
\emph{دوجے درجے دے !!{formula}s} دا سیٹ~$\FrmSOL[L]$،
زبان~$\Lang L$ لئی،
\olref[fol][syn][frm]{defn:formulas} دی تعریف وچ
ایہہ شقّاں شامل کر کے متعین ہوندا اے:
\begin{enumerate}
\item جے $X$ اک $n$-رکنی محمولی متغیر اے تے $t_1$, \dots,
  $t_n$، زبان~$\Lang L$ دے دوجے درجے دے ٹرم نیں، تاں
  $\Atom{X}{t_1,\ldots, t_n}$ اک ایٹمی !!{formula} اے۔

\tagitem{prvAll}{جے $!A$، !!a{formula} اے تے $u$ اک فنکشنی متغیر اے،
  تاں $\lforall[u][!A]$ وی !!a{formula} اے۔}{}

\tagitem{prvAll}{جے $!A$، !!a{formula} اے تے $X$ اک محمولی متغیر اے،
  تاں $\lforall[X][!A]$ وی !!a{formula} اے۔}{}

\tagitem{prvEx}{جے $!A$، !!a{formula} اے تے $u$ اک فنکشنی متغیر اے،
  تاں $\lexists[u][!A]$ وی !!a{formula} اے۔}{}

\tagitem{prvEx}{جے $!A$، !!a{formula} اے تے $X$ اک محمولی متغیر اے،
  تاں $\lexists[X][!A]$ وی !!a{formula} اے۔}{}
\end{enumerate}
\end{defn}

\end{document}
